% S4 certified tetrahedral hinged mechanism --- paper source
%
% Public-facing claims are controlled by paper/CLAIM_BOUNDARY.md and
% docs/PUBLIC_CLAIM_BOUNDARY.md.  RW12 is supplementary digital material only.

\documentclass[11pt,a4paper]{article}


% Reproducible pdfLaTeX output: avoid timestamp/trailer nondeterminism so the
% built PDF can be hashed in the public package manifest.
\ifdefined\pdfinfoomitdate\pdfinfoomitdate=1\fi
\ifdefined\pdfsuppressptexinfo\pdfsuppressptexinfo=-1\fi
\ifdefined\pdftrailerid\pdftrailerid{}\fi

\usepackage[T1]{fontenc}
\usepackage[utf8]{inputenc}
\usepackage{lmodern}
\usepackage{amsmath,amssymb,amsthm,mathtools}
\usepackage{booktabs}
\usepackage{array}
\usepackage{enumitem}
\usepackage{microtype}
\usepackage[margin=1in]{geometry}
\usepackage{xcolor}
\usepackage{xurl}
\usepackage{hyperref}

\hypersetup{
  colorlinks=true,
  linkcolor=blue!55!black,
  citecolor=blue!55!black,
  urlcolor=blue!55!black,
  pdftitle={A Scoped Zero-Thickness S4 Tetrahedral Hinge-Tree Certificate},
  pdfauthor={Oleksiy Babanskyy},
  pdfsubject={S4 tetrahedral hinged mechanism certificate package},
  pdfkeywords={tetrahedron, hinged mechanism, computational certificate, route wrapper, digital fabrication}
}
\emergencystretch=2em

\newtheorem{theorem}{Theorem}[section]
\newtheorem{proposition}[theorem]{Proposition}
\newtheorem{certprop}[theorem]{Certified Proposition}
\newtheorem{lemma}[theorem]{Lemma}
\newtheorem{corollary}[theorem]{Corollary}
\theoremstyle{definition}
\newtheorem{definition}[theorem]{Definition}
\theoremstyle{remark}
\newtheorem{remark}[theorem]{Remark}

\newcommand{\R}{\mathbb R}
\newcommand{\T}{\mathcal T}
\newcommand{\conv}{\operatorname{conv}}
\newcommand{\Vol}{\operatorname{Vol}}
\newcommand{\Rot}{\operatorname{Rot}}
\newcommand{\sgn}{\operatorname{sgn}}
\newcommand{\Pairs}{\mathcal P}
\newcommand{\treeA}{\texttt{TREE\_007}}
\newcommand{\treeB}{\texttt{TREE\_021}}

\title{A Scoped Zero-Thickness S4 Tetrahedral Hinge-Tree Certificate:\\
Median-Plane Geometry, One-Parameter Route Wrappers, and a Selected-Hinge Contact Bridge}
\author{Oleksiy Babanskyy}
\date{}

\begin{document}
\maketitle

\begin{abstract}
We give a scoped, reproducible certificate package for a specific four-piece
(``S4'') dissection of the regular tetrahedron by two opposite-edge median
planes.  The static dissection consists of four congruent tetrahedra whose
closed union is the ambient tetrahedron.  The moving model is a zero-thickness
rigid-piece hinge-tree model on two audited representative signed rays,
\treeA{} and \treeB{}, with parameter domain
\[
  \theta=0 \quad\text{and}\quad 0<\theta\le 120^\circ.
\]
At \(\theta=0\), the pieces form a closed-contact tiling with no strict
interior overlap.  On the open ray, every unordered piece-pair obligation in
the two representative trees is routed to one of three certified predicate
layers: six selected-hinge B04 rows, four common-edge B05 rows, and two
shared-face B06/B07 rows; route-clean B03 obligations are vacuous on this
one-parameter spine.

The main theorem is intentionally narrow.  It is not a proof of physical
hingeability, positive-thickness printability, global collision-free motion, or
three-parameter bounded-cell closure.  The nontrivial semantic point is the
selected-hinge layer: an A7c signed-orientation certificate is not used alone.
A package-local B04/A7c bridge records six row-level wedge certificates showing
that the selected opening side implies boundary hinge contact and no strict
local interior overlap at the contact plane, excluding the hinge axis.  The
repository also includes RW12, a corrected body-preserving digital fabrication
handoff for \treeA{}; RW12 is supplementary digital material and not physical
validation.
\end{abstract}

\noindent\textbf{Keywords.} regular tetrahedron; hinged mechanisms; rigid-piece
kinematics; computational certificates; Sturm certificates; route wrappers;
digital fabrication handoff.

\medskip
\noindent\textbf{Companion extension note.} On the review branch
\texttt{theta-star-addendum-review}, this repository also contains a standalone
addendum-review draft for the equal-magnitude zero-thickness S4 theta-star
finite atlas under \texttt{extensions/theta\_star\_finite\_atlas/}.  That draft
has its own TeX/PDF, bibliography, proof-spine package, and review gates.  It is
not used in the theorem proved below and does not broaden the public claim
boundary of this paper.

\section{Introduction}
\label{sec:introduction}

This paper is a certificate paper for a fixed geometric object, not a general
hinged-dissection theorem.  The object is the median-plane S4 dissection of a
regular tetrahedron.  The four pieces are tetrahedra; at the closed endpoint
they tile the original tetrahedron, and along two audited one-parameter signed
rays they are moved by rooted rigid hinge-tree kinematics.

The certificate is useful precisely because the tempting stronger statements
are false or unproved.  The package does not prove that a positive-thickness
printable mechanism exists.  It does not prove that every possible S4 hinge
choice works.  It does not prove dynamic connectedness between the two audited
representative signed rays.  It proves a scoped statement for two representative
trees and a one-dimensional angular domain, with every pair obligation routed
to a recorded certificate layer.

\paragraph{Contribution.}
The package contributes four pieces of evidence.
\begin{enumerate}[leftmargin=*]
  \item An exact coordinate model for the S4 median-plane dissection and its
  closed-contact endpoint.
  \item A rooted rigid-hinge kinematic model for \treeA{} and \treeB{} with
  explicit hinge axes and signed-ray conventions.
  \item A route-wrapper theorem on the domain
  \(\theta=0\) plus \(0<\theta\le 120^\circ\), supported by the A6--A7d
  certified layers.
  \item A local B04/A7c bridge that upgrades six selected-hinge signed rows
  from mere orientation signs to contact-side/no-strict-local-overlap records.
\end{enumerate}

Replay-backed claims are labelled as \emph{Certified Propositions}.  These are
finite or symbolic certificate statements tied to files in the repository, not
claims of a new general theorem in the literature.

\paragraph{Public boundary.}
The claim boundary is part of the result.  The package allows the zero-thickness
one-parameter route-wrapper claim and the six-row selected-hinge bridge claim.
It disallows physical validation, positive-thickness hinge validity, global
collision-free motion, and three-parameter bounded-cell closure.  These
boundaries are mirrored in \path{docs/PUBLIC_CLAIM_BOUNDARY.md},
\path{docs/CLAIM_LEDGER.md}, and \path{docs/PROOF_OBLIGATIONS.md}.

\paragraph{Organization.}
Section~\ref{sec:model} gives the exact tetrahedral model.  Section~\ref{sec:endpoint}
records the closed endpoint and kinematic conventions.  Section~\ref{sec:wrapper}
states the scoped one-parameter theorem wrapper.  Section~\ref{sec:layers}
explains the A6--A7d certificate layers.  Section~\ref{sec:b04} proves the
selected-hinge wedge bridge used by the B04 rows.  Sections~\ref{sec:package}
and~\ref{sec:rw12} describe reproducibility and the RW12 digital-fabrication
boundary.  Section~\ref{sec:related} places the result next to related hinged
mechanism and fabrication work.

\section{The S4 median-plane model}
\label{sec:model}

Let \(\T=\conv(A,B,C,D)\subset\R^3\) be the unit-edge regular tetrahedron with
\[
A=(0,0,0),\quad B=(1,0,0),\quad
C=\left(\frac{1}{2},\frac{\sqrt3}{2},0\right),\quad
D=\left(\frac{1}{2},\frac{\sqrt3}{6},\frac{\sqrt6}{3}\right).
\]
Let
\[
M_{AB}=\frac{A+B}{2}=\left(\frac{1}{2},0,0\right),\qquad
M_{CD}=\frac{C+D}{2}=\left(\frac{1}{2},\frac{\sqrt3}{3},\frac{\sqrt6}{6}\right).
\]
The S4 pieces are the four closed tetrahedra
\[
\begin{aligned}
P_0&=\conv(A,M_{AB},C,M_{CD}),&
P_1&=\conv(A,M_{AB},D,M_{CD}),\\
P_2&=\conv(B,M_{AB},C,M_{CD}),&
P_3&=\conv(B,M_{AB},D,M_{CD}).
\end{aligned}
\]

\begin{lemma}[Median-plane S4 tiling]
\label{lem:tiling}
The tetrahedron \(\T\) is partitioned by the two barycentric median planes
\(\alpha=\beta\) and \(\gamma=\delta\) into the four closed pieces
\(P_0,P_1,P_2,P_3\).  Their relative interiors are pairwise disjoint, their
closed union is \(\T\), and each piece has volume \(\sqrt2/48\).  The four
pieces are congruent tetrahedra with edge-length multiset
\[
\left\{ \frac{1}{2},\frac{1}{2},\frac{\sqrt2}{2},\frac{\sqrt3}{2},\frac{\sqrt3}{2},1\right\}.
\]
\end{lemma}

\begin{proof}
Write a point of \(\T\) in barycentric coordinates
\[
x=\alpha A+\beta B+\gamma C+\delta D,
\qquad \alpha,\beta,\gamma,\delta\ge0,
\qquad \alpha+\beta+\gamma+\delta=1.
\]
The plane through \(C,D,M_{AB}\) is \(\alpha=\beta\), and the plane through
\(A,B,M_{CD}\) is \(\gamma=\delta\).  For example, on the cell
\(\alpha\ge\beta\), \(\gamma\ge\delta\), one has
\[
x=(\alpha-\beta)A+2\beta M_{AB}+(\gamma-\delta)C+2\delta M_{CD},
\]
with nonnegative coefficients summing to one, hence \(x\in P_0\).  The other
three sign cells give the analogous formulas for \(P_1,P_2,P_3\).  Opposite
strict sign inequalities cannot hold simultaneously, so relative interiors are
disjoint.

The volume calculation
\[
\Vol(P_0)=\frac{1}{6}\left|\det(M_{AB}-A, C-A, M_{CD}-A)\right|=\frac{\sqrt2}{48}
\]
applies the same volume for the other pieces by the ambient symmetries
interchanging \(A,B\) and \(C,D\).  The listed edge lengths follow by direct
coordinate differences for \(P_0\), and symmetry gives the same multiset for
all pieces.
\end{proof}

The closed pair-contact ledger is fixed throughout the package.
\begin{center}
\begin{tabular}{llll}
\toprule
Pair & Contact ID & Contact type & Shared set\\
\midrule
\(P_0-P_1\) & C0 & shared face & \(\conv(A,M_{AB},M_{CD})\)\\
\(P_0-P_2\) & C1 & shared face & \(\conv(C,M_{AB},M_{CD})\)\\
\(P_0-P_3\) & C2 & shared edge & \(\conv(M_{AB},M_{CD})\)\\
\(P_1-P_2\) & C3 & shared edge & \(\conv(M_{AB},M_{CD})\)\\
\(P_1-P_3\) & C4 & shared face & \(\conv(D,M_{AB},M_{CD})\)\\
\(P_2-P_3\) & C5 & shared face & \(\conv(B,M_{AB},M_{CD})\)\\
\bottomrule
\end{tabular}
\end{center}

\begin{certprop}[Source-lock for notation]
The coordinate, piece, pair, contact, and hinge labels in this paper are locked
by \path{certified/source_docs/S4_LEMMA_00_DEFINITIONS_AND_NOTATION_LOCK.md}
and \path{certified/source_docs/S4_LEMMA_01_GEOMETRY_AND_TILING.md}.  Later
certificate rows are interpreted using these labels.
\end{certprop}

\section{Closed endpoint and rooted hinge-tree kinematics}
\label{sec:endpoint}

The endpoint \(\theta=0\) is a closed-contact endpoint, not a positive-clearance
endpoint.  The pieces are closed sets and are allowed to meet along the
catalogued contacts in Section~\ref{sec:model}.

\begin{lemma}[Closed-contact endpoint]
\label{lem:endpoint}
At \(\theta=0\), the four pieces lie in \(\T\), their volumes sum to
\(\Vol(\T)\), all six unordered pairs are exactly the catalogued contacts C0--C5,
and no two distinct pieces have strict interior overlap.  This endpoint has
zero clearance at the catalogued contacts.
\end{lemma}

\begin{proof}
By Lemma~\ref{lem:tiling}, the closed union is \(\T\) and relative interiors
are pairwise disjoint.  Thus every pairwise intersection is contained in a
proper face or edge of each piece and cannot contain a three-dimensional
interior point common to two pieces.  The contact table in Section~\ref{sec:model}
lists all six unordered piece pairs.  Since the list includes shared faces and
shared edges, the endpoint is contact-allowing and not positive-clearance.
\end{proof}

The moving model uses rooted hinge trees.  For an oriented axis from \(a\) to
\(b\), write \(u=(b-a)/\lVert b-a\rVert\).  The rotation about the affine line
through \(a,b\) by angle \(\phi\) is
\[
\Rot(a,b,\phi)(x)=Q_u(\phi)x+a-Q_u(\phi)a,
\]
where \(Q_u\) is Rodrigues' rotation matrix.  Points on the axis are fixed.

The audited representative trees are:
\begin{center}
\begin{tabular}{lll}
\toprule
Tree & Hinge graph & Signed hinge angles\\
\midrule
\treeA{} & \(P_0P_1, P_0P_2, P_1P_3\) &
\(d_{H0}=\theta,\ d_{H4}=\theta,\ d_{H7}=-\theta\)\\
\treeB{} & \(P_0P_1, P_1P_3, P_2P_3\) &
\(d_{H0}=\theta,\ d_{H7}=-\theta,\ d_{H9}=\theta\)\\
\bottomrule
\end{tabular}
\end{center}
The root is \(P_0\) in both cases.

\begin{lemma}[Well-defined rooted kinematics]
\label{lem:kinematics}
For \treeA{} and \treeB{}, any real assignment of signed angles to the selected
hinges determines a unique rigid transform \(F_i\) for each piece \(P_i\).  For
each selected hinge edge, the parent and child transforms agree on the hinge
axis.
\end{lemma}

\begin{proof}
Each selected graph has four vertices, three edges, and the rooted paths listed
in the source lock; hence it is a tree.  There is a unique path from root
\(P_0\) to every piece.  Starting with \(F_0=\mathrm{id}\), propagate along each
edge by composing the parent transform with the axis rotation about the parent-
placed hinge line.  A composition of rigid transforms is rigid, and uniqueness
of the rooted path gives uniqueness of \(F_i\).  Since \(\Rot(a,b,\phi)\) fixes
the axis line pointwise, parent and child transforms agree on the selected
hinge axis.
\end{proof}

At \(\theta=0\), every signed angle is zero, every axis rotation is the identity,
and all piece transforms are identity.  Thus both representative trees have the
same closed-contact endpoint from Lemma~\ref{lem:endpoint}.

\begin{certprop}[Kinematic source-lock]
The endpoint and kinematic statements are source-locked to
\path{certified/source_docs/S4_LEMMA_02_CLOSED_ENDPOINT.md} and
\path{certified/source_docs/S4_LEMMA_03_KINEMATICS_AND_SIGNS.md}.  These files
also record the exact hinge IDs \texttt{H0\_A\_M\_AB}, \texttt{H4\_C\_M\_CD},
\texttt{H7\_D\_M\_CD}, and \texttt{H9\_B\_M\_AB}.
\end{certprop}

\section{The scoped one-parameter wrapper}
\label{sec:wrapper}

The one-parameter open domain is
\[
0<\theta\le120^\circ.
\]
The certificate uses the half-angle coordinate
\[
t=\tan(\theta/2),\qquad 0<t\le\sqrt3,
\]
and many sign proofs are carried out on the rational open superset
\[
0<t<2.
\]

Every representative tree has six unordered pair obligations.  The A7d wrapper
routes the twelve obligations across the two trees as follows.
\begin{center}
\begin{tabular}{lll}
\toprule
Tree & Pair & Route\\
\midrule
\treeA{} & \(P_0-P_1\) & B04 selected-hinge contact-side\\
\treeA{} & \(P_0-P_2\) & B04 selected-hinge contact-side\\
\treeA{} & \(P_0-P_3\) & B05 common-edge projection\\
\treeA{} & \(P_1-P_2\) & B05 common-edge projection\\
\treeA{} & \(P_1-P_3\) & B04 selected-hinge contact-side\\
\treeA{} & \(P_2-P_3\) & B06/B07 shared-face residual\\
\midrule
\treeB{} & \(P_0-P_1\) & B04 selected-hinge contact-side\\
\treeB{} & \(P_0-P_2\) & B06/B07 shared-face residual\\
\treeB{} & \(P_0-P_3\) & B05 common-edge projection\\
\treeB{} & \(P_1-P_2\) & B05 common-edge projection\\
\treeB{} & \(P_1-P_3\) & B04 selected-hinge contact-side\\
\treeB{} & \(P_2-P_3\) & B04 selected-hinge contact-side\\
\bottomrule
\end{tabular}
\end{center}

\begin{theorem}[Scoped zero-thickness one-parameter wrapper]
\label{thm:wrapper}
For the catalogued median-plane S4 pieces, the public package records a closed
zero-thickness one-parameter wrapper for the two audited representatives
\treeA{} and \treeB{} on the domain
\[
\theta=0\quad\text{and}\quad 0<\theta\le120^\circ.
\]
At \(\theta=0\), the configuration is the closed-contact endpoint of
Lemma~\ref{lem:endpoint}.  On the open ray, every unordered pair obligation in
each representative tree is assigned to a completed certificate layer: six B04
selected-hinge bridge rows, four B05 common-edge projection rows, and two
B06/B07 shared-face residual rows.  There are no route-clean B03 obligations on
this one-parameter spine.
\end{theorem}

\begin{proof}
The endpoint clause follows from Lemmas~\ref{lem:tiling}--\ref{lem:kinematics}:
at \(\theta=0\), all rooted transforms are identity and the pieces are in the
catalogued closed tiling.

For the open ray, the A7d manifest is a pair-route ledger.  It contains two
wrapper records, one for \treeA{} and one for \treeB{}, each with six unordered
piece-pair obligations.  The manifest records the predicate route counts
\[
\text{B04}=6,
\qquad
\text{B05}=4,
\qquad
\text{B06/B07}=2.
\]
The B03 count is zero by the A7b route-vacuity certificate.  The routed layers
are discharged by the certified propositions in Sections~\ref{sec:layers} and
\ref{sec:b04}.  Thus every open-ray pair obligation is routed to a completed
package-local layer, and no additional B03 obligation remains on this scoped
one-parameter domain.
\end{proof}

\begin{certprop}[A7d manifest]
The theorem wrapper above is the statement recorded in
\path{certified/a7d_one_parameter_theorem_wrapper_manifest.json}.  The manifest
has \texttt{record\_count=2}, \texttt{wrapper\_closed\_count=2}, and object
status count
\texttt{a7d\_one\_parameter\_ray\_wrapper\_closed: 2}.
\end{certprop}

\begin{remark}[Scope]
The theorem is a route-wrapper theorem for a zero-thickness one-parameter model.
It is not a positive-thickness fabrication theorem, not a global collision-free
motion theorem, and not a three-parameter bounded-cell theorem.
\end{remark}

\section{Certificate layers A6--A7d}
\label{sec:layers}

This section records the certificate layers used by Theorem~\ref{thm:wrapper}.
The layers are package-local and source-backed; their role is to discharge the
routed pair obligations, not to prove a broader mechanism theorem.

\begin{certprop}[A6: B05 common-edge projection]
The B05 common-edge projection layer is symbolically closed on the one-parameter
ray for the relevant records.  A6 packages the A3/A4 and A5 algebraic
certificates: A3/A4 proves full-domain raw gap, normalized gap, and axis
positivity on the open ray, while A5 splits at \(t=1\) and proves post-switch
support signatures and projection gaps on \(1<t<7/4\), which contains
\(1<t\le\sqrt3\).  The A7d wrapper uses four B05 pair routes.
\end{certprop}

\begin{proof}[Certification]
The source is
\path{certified/source_docs/S4_CL5_A6_ONE_PARAMETER_RAY_CLOSURE_PACKAGE.md}.
It records seven symbolic B05 closures in the source package and the one-
parameter B05 closure chain used by the wrapper.  The public wrapper consumes
the B05 role through \path{certified/a7d_one_parameter_theorem_wrapper_manifest.json}.
\end{proof}

\begin{certprop}[A7a: B06/B07 shared-face residual]
For the two one-parameter shared-face residual targets,
\[
\treeA{}:P_2-P_3,
\qquad
\treeB{}:P_0-P_2,
\]
the residual shared-face formula is positive on the open ray.  In half-angle
coordinate the certified rational form is
\[
\frac{t^3}{(1+t^2)^2},
\qquad 0<t<2.
\]
In particular it is positive on \(0<t\le\sqrt3\).
\end{certprop}

\begin{proof}
The trigonometric formula from the source layer is
\[
\frac{(1-\cos\theta)\sin\theta}{4}
  =\sin^3(\theta/2)\cos(\theta/2).
\]
With \(t=\tan(\theta/2)\), this becomes
\[
\frac{t^3}{(1+t^2)^2}.
\]
The numerator and denominator have no open roots on \(0<t<2\); at \(t=1\) both
are positive.  Therefore the expression is positive on the certified open
superset, hence on the audited ray domain.  The Sturm/root-count certificate is
recorded in
\path{certified/source_docs/S4_CL5_A7A_SHARED_FACE_RESIDUAL_STURM_CERTIFICATE.md}.
\end{proof}

\begin{certprop}[A7b: B03 route vacuity]
There are no route-clean ordinary non-contact B03 obligations on the audited
one-parameter representative rays.  Across the two representatives, the twelve
pair obligations route as six selected-hinge contacts, four residual shared-edge
common-edge projections, and two residual shared-face targets.
\end{certprop}

\begin{proof}[Certification]
This is the A7b route-vacuity certificate in
\path{certified/source_docs/S4_CL5_A7B_B03_RAY_VACUITY_CERTIFICATE.md}.  It
records \texttt{one-parameter pair records: 12}, \texttt{B03 route-clean pairs:
0}, and \texttt{B03 Sturm obligations: 0}.
\end{proof}

\begin{certprop}[A7c: selected-hinge signed orientation]
For the six selected B04 rows, A7c certifies the sign of the selected hinge
angle on the open ray.  The signed orientation has the same sign as \(s_h t\)
on \(0<t<2\), with ray-sign counts \(+1:4\) and \(-1:2\).
\end{certprop}

\begin{proof}
Lemma~\ref{lem:kinematics} fixes the signed-ray convention
\(d_h(\theta)=s_h\theta\).  In half-angle coordinates, the relevant signed
orientation is sign-equivalent to \(s_h t\).  The only root is at \(t=0\),
which is excluded from the open ray.  The sample point \(t=1\) fixes the sign.
The six rows are recorded in
\path{certified/a7c_selected_hinge_contact_side_certificate_manifest.json} and
\path{certified/source_docs/S4_CL5_A7C_SELECTED_HINGE_CONTACT_SIDE_CERTIFICATE.md}.
\end{proof}

A7c alone is not the B04 contact-side proof.  It is the signed-orientation
source layer consumed by the bridge in Section~\ref{sec:b04}.

\section{The B04/A7c selected-hinge bridge}
\label{sec:b04}

The B04 rows are zero-margin selected hinge contacts.  They cannot be certified
by positive SAT clearance: the pieces are meant to remain joined along the
hinge axis.  The correct local statement is boundary contact plus no strict
local interior overlap away from the hinge axis.

Fix one selected row.  Let the oriented hinge axis be \(a\to b\), and let
\(F\) be the third vertex of the source shared face, so the source face is
\(\conv(a,b,F)\).  Let \(Q_p\) be the parent apex not in this face and \(Q_c\)
the child apex not in this face.  Define the oriented side form
\[
\Omega(X,Y)=(b-a)\cdot((X-a)\times(Y-a)).
\]

\begin{lemma}[Local selected-hinge wedge]
\label{lem:wedge}
For each of the six selected B04 rows, with A7c ray sign \(\sigma\), the package
records the exact side identities
\[
\sigma\Omega(F,Q_p)=-\frac{\sqrt2}{8}<0,
\qquad
\sigma\Omega(F,Q_c)=\frac{\sqrt2}{8}>0.
\]
Thus the A7c sign identifies the child opening side of the selected source
face.  Moreover the cross-section angle \(\alpha\) between the source face ray
and either adjacent tetrahedral apex ray satisfies
\[
\cos\alpha=\frac{\sqrt6}{3}>\frac{1}{2},
\qquad\text{so}\qquad
0<\alpha<60^\circ.
\]
Consequently, for \(0<\theta\le120^\circ\), the opened child sector satisfies
\(\theta+\alpha<180^\circ\) and cannot wrap behind the hinge axis into the
parent sector.
\end{lemma}

\begin{proof}
The identities for \(\Omega\) are exact coordinate evaluations in the normal
cross-section about the selected hinge axis.  They place the parent apex and
child apex on opposite sides of the oriented source face, after multiplying by
the row sign \(\sigma\).  The angle calculation gives
\(\cos\alpha=\sqrt6/3>1/2\), hence \(\alpha<60^\circ\).  Since the ray domain
has \(\theta\le120^\circ\), the sector opened by the child face remains within
an open half-turn from the source face.  Therefore it cannot cross through the
opposite parent sector away from the hinge axis.
\end{proof}

The bridge also checks the actual rotated child vertices using
\(t=\tan(\theta/2)\).  For the rotated third vertex of the selected source face,
the signed side is
\[
\frac{t}{2(t^2+1)}>0,
\qquad 0<t\le\sqrt3.
\]
For the rotated child apex, each row reduces to a positive rational form
recorded in the row JSON: the denominator is positive, the sample signs at
\(t=1\) and \(t=\sqrt3\) are positive, and all positive numerator roots lie
strictly above \(\sqrt3\).

\begin{certprop}[B04/A7c bridge rows]
\label{prop:b04-bridge}
The package-local bridge accepts exactly six records:
\begin{center}
\begin{tabular}{llllr}
\toprule
Tree & Pair & Hinge & Contact & Ray sign\\
\midrule
\treeA{} & \(P_0-P_1\) & H0\_A\_M\_AB & C0 & \(+1\)\\
\treeA{} & \(P_0-P_2\) & H4\_C\_M\_CD & C1 & \(+1\)\\
\treeA{} & \(P_1-P_3\) & H7\_D\_M\_CD & C4 & \(-1\)\\
\treeB{} & \(P_0-P_1\) & H0\_A\_M\_AB & C0 & \(+1\)\\
\treeB{} & \(P_1-P_3\) & H7\_D\_M\_CD & C4 & \(-1\)\\
\treeB{} & \(P_2-P_3\) & H9\_B\_M\_AB & C5 & \(+1\)\\
\bottomrule
\end{tabular}
\end{center}
For these rows, A7c signed orientation plus the local wedge bridge proves
boundary hinge contact and no strict local interior overlap at the contact
plane, excluding the hinge axis.
\end{certprop}

\begin{proof}[Certification]
The bridge manifest is
\path{certified/b04_a7c_contact_side_bridge_manifest.json}.  It records
\texttt{record\_count=6}, \texttt{accepted\_bridge\_record\_count=6}, and the
policy that bridge rows are consumed from the A7c manifest records while
asserting equality of \texttt{tree\_id}, \texttt{hinge\_id}, \texttt{pair}, and
\texttt{ray\_sign}.  The row records live in
\path{certified/a7c_b04_contact_side_bridge_records/}.  The derivation and
claim boundary are source-locked in
\path{certified/source_docs/S4_CL5_A7C_B04_CONTACT_SIDE_BRIDGE_CERTIFICATE.md}.
\end{proof}

\begin{remark}
The bridge is local and selected-row scoped.  It does not give positive
separation for selected hinges, does not handle hinge thickness, and does not
extend to unselected or three-parameter B04 cases.
\end{remark}

\section{Public package and reproducibility}
\label{sec:package}

The repository is organized as a paper-as-public-package: the manuscript,
certificate manifests, row-level records, replay scripts, hash manifest, and
claim-boundary ledgers are shipped together.  This follows established
reproducible-artifact practice rather than asserting a new standard; its role
here is to make the computational dependencies auditable while separating the
mathematical theorem from supplementary digital-fabrication artifacts
\cite{smith2017joss,rougier2017rescience}.  The main paper source and PDF live
in \path{paper/}.  Mathematical and certificate controls live in
\path{docs/}, \path{certified/}, \path{scripts/}, and \path{results/}.  The
public package manifest is \path{paper/PUBLIC_PACKAGE_MANIFEST.json}.

\begin{certprop}[Public package gate]
The local public package gate passes.  The command
\begin{center}
\path{python scripts/run_all_reproducibility_checks.py}
\end{center}
checks required files, B04/A7c bridge row counts, RW12 hash consistency,
manifest hashes, and absence of promoted physical-validation phrases.  The
current public package check records all checks passing.
\end{certprop}

\begin{proof}[Certification]
The generated report is \path{results/public_package_check.json}.  The pytest
wrapper \path{tests/test_public_package.py} runs the same package checker.  The
manifest records \texttt{LOCAL\_RELEASE\_GATES\_PASSED},
\texttt{public\_export\_ready=true}, and
\texttt{physical\_validation\_claimed=false}.
\end{proof}

The source-lock boundary is intentionally duplicated: the paper cites concise
Certified Propositions, while the repository ships the longer source documents
under \path{certified/source_docs/}.  In particular, the public package includes
S4 Lemmas 00--12 and the B04 contact-orientation review notes, so the bridge
source dependencies are present in the package rather than only in the private
workspace.

\begin{remark}[What the tests do not prove]
The public tests verify package integrity and row-level consistency.  They do
not replace the mathematical content of Lemmas~\ref{lem:tiling}--\ref{lem:wedge}
and Theorem~\ref{thm:wrapper}; they keep the shipped artifacts synchronized with
the claim boundary.
\end{remark}

\section{RW12 supplementary digital fabrication handoff}
\label{sec:rw12}

RW12 is included because it is practically useful, not because it strengthens
Theorem~\ref{thm:wrapper}.  It is a corrected body-preserving digital handoff
for \treeA{} containing fabrication-oriented files, checklist material, and a
compact package zip.  Its governing report is
\path{results/rw12_external_fabrication_review_package/rw12_external_fabrication_review_report.json}.

\begin{certprop}[RW12 hash and claim boundary]
The public checker verifies that the RW12 zip hash matches the RW12 report and
that RW12 does not promote a physical-validation claim.  The RW12 package is
therefore an appendix-level digital handoff, not theorem evidence.
\end{certprop}

The practical interpretation is straightforward: RW12 can be passed to a CAD,
fabrication, or external-review workflow as a digital artifact.  It does not
certify positive-thickness hinge clearances, material tolerances, assembly
success, printed prototype behavior, or measured motion.

\begin{remark}
This distinction matters for publication.  The mathematical result is the
zero-thickness one-parameter certificate package.  A real-world mechanism would
require additional positive-thickness modelling, tolerances, fabrication
measurements, and physical validation.  Those are future work, not hidden
assumptions of the present theorem.
\end{remark}

\section{Related work and remaining obligations}
\label{sec:related}

The general hinged-dissection landscape is much broader than this paper.  Hinged
dissections exist in substantial generality~\cite{abbott2007hinged}; reversible
hinged dissections and reversible nets are studied in
\cite{akiyama2018polyhedral,akiyama2016reversible}; and standard background on
geometric folding, linkages, and polyhedra is given by Demaine and
O'Rourke~\cite{demaine2007gfalop}.  Continuous flattening, vertex-unfolding,
and digital jointed assembly systems provide adjacent contexts
\cite{demaine2024continuous,demaine2001vertex,rodsteward2019,joinable2021}.

The novelty candidate here is not a new universal hinged-dissection theorem.
It is the combination of a fixed tetrahedral median-plane model, exact
zero-thickness rooted hinge-tree kinematics, row-level route-wrapper
certificates, and a public companion repository that keeps theorem evidence separate
from digital fabrication artifacts.

\paragraph{Remaining obligations.}
The following are deliberately open or out of scope.
\begin{itemize}[leftmargin=*]
  \item Positive-thickness hinge design and physical fabrication validation.
  \item Global collision-free motion beyond the audited one-parameter domain.
  \item Dynamic connectedness between \treeA{} and \treeB{}.
  \item Three-parameter bounded-cell closure.
  \item Extension of the B04 bridge beyond the six selected rows.
  \item A general theorem for all S4 hinge trees or all tetrahedral dissections.
\end{itemize}

The package is therefore best read as a scoped certificate paper: it identifies
what has been proved for a concrete object, ships the evidence in a reproducible
layout, and makes the boundaries visible enough for independent review.


\bibliographystyle{plain}
\bibliography{refs}

\end{document}
